% id: 124-16
% book: 베이직쎈 확률과 통계 · 07 통계적 추정
% section: 개념 38 표본평균의 평균, 분산, 표준편차
% figure: none
% answer: 
% answer_source: 
% variant_level: 0 (원본 전사)
% 이 파일은 build.mjs 가 items.json 에서 만든다. 고칠 때는 items.json 을 고친다.
\smash{\raisebox{1.5mm}{\dmkichul{베이직쎈 확률과 통계 124쪽 16번}}}
모집단 0, 2, 4에서 크기가 2인 표본을 임의추출할 때, 다음을 구하시오.
\dmsubprob{1} 표본평균 $\overline{X}$가 가질 수 있는 값\dmhint{표본이 $(0{,}\mkern3mu\mathopen{}0)$일 때,\\ $\overline{X}=\dfrac{0+0}{2}=0$\\ 표본이 $(0{,}\mkern3mu\mathopen{}2)$, $(2{,}\mkern3mu\mathopen{}0)$일 때,\\ $\overline{X}=\dfrac{0+2}{2}=1$\\ 표본이 $(0{,}\mkern3mu\mathopen{}4)$, $(4{,}\mkern3mu\mathopen{}0)$, $(2{,}\mkern3mu\mathopen{}\blank)$일 때,\\ $\overline{X}=\dfrac{0+\blank}{2}=\dfrac{2+\blank}{2}=\blank$\\ 표본이 $(2{,}\mkern3mu\mathopen{}\blank)$, $(\blank{,}\mkern3mu\mathopen{}2)$일 때,\\ $\overline{X}=\dfrac{2+\blank}{2}=\blank$\\ 표본이 $(4{,}\mkern3mu\mathopen{}4)$일 때,\\ $\overline{X}=\dfrac{4+4}{2}=4$\\ 따라서 $\overline{X}$가 가질 수 있는 값은\\ $0{,}\mkern3mu\mathopen{}1{,}\mkern3mu\mathopen{}\blank{,}\mkern3mu\mathopen{}\blank{,}\mkern3mu\mathopen{}4$}
\dmsubprob{2} 표본평균 $\overline{X}$의 확률분포를 나타낸 표\par\smallskip\centerline{{\setlength{\tabcolsep}{5pt}\renewcommand{\arraystretch}{2.2}\begin{tabular}{|c|c|c|c|c|c|c|}\hline $\overline{X}$ & $0$ & $1$ & \makebox[2.2em]{} & \makebox[2.2em]{} & $4$ & 합계 \\ \hline $\mathrm{P}(\overline{X}=\overline{x})$ & $\dfrac{1}{9}$ & & & $\dfrac{2}{9}$ & & \\ \hline\end{tabular}}}\par\smallskip
\dmsubprob{3} $\mathrm{E}(\overline{X})$, $\mathrm{V}(\overline{X})$, $\sigma(\overline{X})$
